\documentclass[11pt,a4paper]{article}

%% PACHETE MATEMATICE
\usepackage{amsmath,amssymb,amsfonts}
\usepackage{amsthm}
\usepackage{mathpazo}
\usepackage{eulervm}
\usepackage{enumerate}
\usepackage{color}
\usepackage{graphicx}
\usepackage[all]{xy}
\usepackage{xcolor}
\usepackage{framed}
\usepackage[unicode]{hyperref}
\setcounter{MaxMatrixCols}{20}
\usepackage{multicol}
%\usepackage[romanian]{babel}


%% ASPECT
\linespread{1.05}
\usepackage[T1]{fontenc}
\usepackage[utf8x]{inputenc}
\usepackage[margin=2cm]{geometry}
% \usepackage[color]{showkeys}
\usepackage{anyfontsize}
\usepackage{titlesec}
\titleformat*{\section}{\large\bfseries}

\usepackage{fancyhdr}
\pagestyle{fancy}
\rhead{\small \color{gray} Notițe de Adrian Manea}
\lhead{\small \color{gray} Aero-BCD, 2017---2018, Prof. L. Costache \& M. Olteanu}


\usepackage[autostyle = false, style = german]{csquotes}
\usepackage[romanian]{babel}
\newcommand\qq{\enquote}

%% COMENZILE MELE
\renewcommand*{\proofname}{Dem.:}
\newcommand\bb{\mathbb}
\newcommand\dr{\mathrm}
\newcommand\kal{\mathcal}
\newcommand\fk{\mathfrak}
\newcommand\op{\oplus}
\newcommand\ot{\otimes}
\newcommand\ds{\displaystyle}
\newcommand\id{\indent}
\newcommand\nid{\noindent}
\newcommand\seq{\subseteq}
\newcommand\sneq{\subsetneq}
\newcommand\speq{\supseteq}
\newcommand\spneq{\supsetneq}
\newcommand\rar{\rightarrow}
\newcommand\Rar{\Rightarrow}
\newcommand\xrar{\xrightarrow}
\newcommand\lar{\leftarrow}
\newcommand\Lar{\Leftarrow}
\newcommand\xlar{\xleftarrow}
\newcommand\lrar{\leftrightarrow}
\newcommand\Lrar{\Leftrightarrow}
\newcommand\ideal{\trianglelefteq}
\newcommand\ai{\text{ a.\^{i}. }}
\newcommand\sk{(\textit{Schi\c{t}\u{a}}) }
\newcommand\rhup{\rightharpoonup}
\newcommand\rhdn{\rightharpoondown}
\newcommand\lhup{\leftharpoonup}
\newcommand\lhdn{\leftharpoondown}
\newcommand\vid{\emptyset}
\newcommand\wh{\widehat}
\newcommand\ol{\overline}
\newcommand\NN{\mathbb{N}}
\newcommand\RR{\mathbb{R}}
\newcommand\ZZ{\mathbb{Z}}
\newcommand\QQ{\mathbb{Q}}
\newcommand\CC{\mathbb{C}}

%% STILURILE MELE
\newtheoremstyle{dotless}{}{}{\itshape}{}{\bfseries}{:}{ }{}

\theoremstyle{dotless}
\newtheorem{theorem}{Teorem\u{a}}[section]
\newtheorem{proposition}{Propozi\c{t}ie}[section]
\newtheorem{lemma}{Lem\u{a}}[section]
\newtheorem{corollary}{Corolar}[section]
\newtheorem{exercise}{Exerci\c{t}iu}

\newtheoremstyle{dotlessdr}{}{}{}{}{\bfseries}{:}{ }{}
\theoremstyle{dotlessdr}
\newtheorem{example}{Exemplu}[section]
\newtheorem{definition}{Defini\c{t}ie}[section]
\newtheorem{remark}{Observa\c{t}ie}[section]




\begin{document}
% \definecolor{labelkey}{RGB}{222,0,0}

\begin{center}
  {\large\textbf{Seminar 5 \\ Spații metrice. Serii de funcții}}
\end{center}

\section{Spații metrice}
\label{sec:sp-metrice}

\begin{definition}\label{def:distanta}
  Fie $X$ o mulțime nevidă. O aplicație $d : X \times X \to \RR$ se numește \emph{distanță}
  (metrică) pe $X$ dacă:
  \begin{enumerate}[(a)]
  \item $d(x, y) \geq 0, \forall x, y \in \RR$;
  \item $d(x, y) = 0 \Leftrightarrow x = y$;
  \item $d(x, y) = d(y, x), \forall x, y \in X$;
  \item $d(x, z) \leq d(x, y) + d(y, z)$ (\emph{inegalitatea triunghiului}).
  \end{enumerate}

  În acest context, perechea $(X, d)$ se numește \emph{spațiu metric}.
\end{definition}

Această noțiune vine să generalizeze calculul distanțelor cu ajutorul modulului,
cum se procedează în cazul mulțimii numerelor reale, de exemplu. În consecință,
avem și următoarele \emph{mulțimi remarcabile în spații metrice $(X, d)$}:
\begin{itemize}
\item \emph{bila deschisă de centru $a$ și rază $r$}, definită pentru $a, r \in \RR$ prin:
  \[
    B(a, r) = \{ x \in X \mid d(a, x) < r \};
  \]
\item \emph{sfera de centru $a$ și rază $r$}, definită prin:
  \[
    S(a, r) = \{ x \in X \mid d(a, x) = r \};
  \]
\item \emph{bila închisă de centru $a$ și rază $r$}, definită prin:
  \[
    \overline{B(a, r)} = B(a, r) \cup S(a, r) = \{x \in X \mid d(a, x) \leq r \}.
  \]
\end{itemize}

Aceste mulțimi generalizează \emph{intervalele de numere reale}. Mai general, avem:
\begin{definition}\label{def:m-deschisa}
  Fie $X$ un spațiu metric. O submulțime $D \seq X$ se numește \emph{deschisă} dacă
  $\forall a \in D, \exists r > 0$ astfel încît $B(a, r) \seq D$. Cu alte cuvinte,
  mulțimea conține o bilă deschisă, centrată în orice punct al său.

  O mulțime se numește \emph{închisă} dacă complementara ei, relativ la spațiul total,
  este o mulțime deschisă.
\end{definition}

Adaptînd noțiunile specifice pentru analiza matematică, precum convergență,
limită ș.a.m.d. folosindu-ne de funcția distanță, putem defini construcțiile
și conceptele uzuale. În particular, au sens noțiuni precum \emph{șiruri convergente}
și \emph{șiruri Cauchy}, definite ca în cazul mulțimii numerelor reale, doar
cu ajutorul funcției generale distanță. În plus, avem și următoarea noțiune:

\begin{definition}\label{def:spatiu-complet}
  Un spațiu metric $(X, d)$ se numește \emph{complet} dacă noțiunile de
  \qq{șir Cauchy} și \qq{șir convergent} coincid pentru $X$.
\end{definition}

%%%%%%%%%%%%%%%%%%%%%%%%%%%%%%%%%%%%%%%%%%%%%%%%%%%%%%%%%%%%%%%%%%%%%%

\section{Principiul contracției. Metoda aproximațiilor succesive}
\label{sec:contractie}

\begin{definition}\label{def:contractie}
  Fie $(X, d)$ un spațiu metric și fie $f : X \to X$ o funcție.

  Aplicația $f$ se numește \emph{contracție pe $X$} dacă există $k \in [0, 1)$
  astfel încît:
  \[
    d(f(x), f(y)) \leq k \cdot d(x, y), \forall x, y \in X.
  \]

  Numărul $k$ se numește \emph{factor de contracție}.
\end{definition}

Un rezultat important este următorul:

\begin{theorem}[Banach]\label{thm:banach-fix}
  Fie $(X, d)$ un spațiu metric complet și fie $f : X \to X$ o contracție de factor
  $k$. Atunci există un unic punct $\xi \in X$, astfel încît $f(\xi) = \xi$.

  În acest context, $\xi$ se numește \emph{punct fix} pentru $f$.
\end{theorem}

Pentru a găsi punctul fix al unei aplicații, se folosește \emph{metoda %
  aproximațiilor succesive}. Se construiește un șir recurent astfel.

Fie $x_0 \in X$, arbitrar și definim șirul recurent $x_{n+1} = f(x_n)$.
Se poate demonstra că șirul $x_n$ este convergent, iar limita sa este punctul
fix căutat. În plus, eroarea aproximației cu acest șir este dată de:
\[
  d(x_n, \xi) \leq \frac{k^n}{1 - k} \cdot d(x_0, x_1), \forall n \in \NN.
\]

%%%%%%%%%%%%%%%%%%%%%%%%%%%%%%%%%%%%%%%%%%%%%%%%%%%%%%%%%%%%%%%%%%%%%%

\section{Exemplu rezolvat}
\label{sec:ex-rez}

Aplicațiile interesante pentru această temă sînt date de calculul aproximativ
al soluțiilor unor ecuații, definite în spații metrice.

1. Să se aproximeze cu o eroare mai mică decît $10^{-3}$ soluția reală a
ecuației $x^3 + 4x - 1 = 0$.

\emph{Soluție:} Folosind, eventual, metode de analiză (șirul lui Rolle),
se poate arăta că ecuația are o singură soluție reală $\xi \in (0, 1)$.
Folosim metoda aproximațiilor succesive pentru a o găsi.

Fie $X = [0, 1]$ și $f: X \to X, f(x) = \frac{1}{x^2 + 4}$. Pînă la o
translație, ecuația dată este echivalentă cu $f(x) = x$, adică a găsi
un punct fix pentru funcția $f$.

Spațiul metric $X$ este complet. Mai demonstrăm că $f$ este contracție
pe $X$. Derivata este $f'(x) = \frac{-2x}{(x^2 + 4)^2}$, și avem:
\[
  \sup_{x \in X} | f'(x) | = -f'(1) = \frac{2}{25} < 1,
\]
deci $f$ este o contracție, de factor $k = \frac{2}{25}$.

Șirul aproximațiilor succesive este dat de:
\[
  x_0 = 0, \quad x_{n + 1} = f(x_n) = \frac{1}{x_n^2 + 4}.
\]

Evaluarea erorii:
\[
  |x_n - \xi| < \frac{k^n}{1 - k}|x_0 - x_1| = \frac{1}{3} \Big( \frac{2}{25} \Big)^n,
\]
de unde rezultă că putem lua $\xi \simeq x_3 = f \big(\frac{16}{65}\big) \simeq 0,235$.

\begin{remark}\label{rk:alternativ}
  Alternativ, puteam lucra cu $g(x) = \frac{1}{4} (1 - x^3)$, cu $x \in [0, 1]$.
  Se arată că și $g$ este o contracție, de factor $k = \frac{3}{4}$. În acest caz,
  șirul aproximațiilor succesive converge mai încet și avem $\xi \simeq x_6$.
\end{remark}

\vspace{1cm}

Similar, puteți rezolva următoarele ecuații, cu eroarea $\varepsilon$:
\begin{enumerate}[(a)]
\item $x^3 + 12x - 1 = 0, \quad \varepsilon = 10^{-3}$;
\item $x^5 + x - 15 = 0, \quad \varepsilon = 10^{-3}$;
\item $3x + e^{-x} = 1, \quad \varepsilon = 10^{-3}$;
\item $x^3 - x + 5 = 0, \quad \varepsilon = 10^{-2}$;
\item $x^5 + 3x - 2 = 0, \quad \varepsilon = 10^{-3}$.
\end{enumerate}

%%%%%%%%%%%%%%%%%%%%%%%%%%%%%%%%%%%%%%%%%%%%%%%%%%%%%%%%%%%%%%%%%%%%%%

\section{Șiruri de funcții}
\label{sec:siruri-f}

\begin{definition}\label{def:pc-uc}
  Fie $(f_n)_n$ un șir de funcții, cu fiecare $f_n : [a, b] \to \RR$ și fie
  o funcție $f: [a, b] \to \RR$. Spunem că șirul $(f_n)$ \emph{converge punctual}
  pe $[a, b]$ către $f$ pentru $n \to \infty$, scris $f_n \xrar{PC} f$ dacă avem
  $f_n(x_0) \to f(x_0)$, pentru orice $x \in [a, b]$.

  Spunem că șirul $(f_n)$ \emph{converge uniform} pe $[a, b]$ către $f$, scris
  $f_n \xrar{UC} f$, dacă:
  \[
    \forall \varepsilon > 0, \exists N_\varepsilon \in \NN \ai %
    \forall n \geq N_\varepsilon, |f_n(x) - f(x)| < \varepsilon, \forall x \in [a, b].
  \]
\end{definition}

În calcule, va fi de folos următorul rezultat:

\begin{theorem}\label{thm:pc-uc}
  Un șir $(f_n)$ de funcții mărginite $f_n : [a, b] \to \RR$ este uniform convergent
  către o funcție $f$ dacă și numai dacă $\ds\lim_{n \to \infty} ||f_n - f|| = 0$.
\end{theorem}

De asemenea, avem și:
\begin{theorem}\label{thm:pc-uc-implica}
  Orice șir de funcții $f_n : [a, b] \to \RR$ uniform convergent pe $[a, b]$ este
  punctual convergent pe $[a, b]$.
\end{theorem}

Reciproca acestui rezultat este falsă: fie $[a, b] = [0, 1]$ și $f_n(x) = x^n, n \geq 1$.

Pentru orice $x \in [a, b]$, avem:
\[
  \lim_{n \to \infty} f_n(x) =
  \begin{cases}
    0, & x \in [0, 1) \\
    1, & x = 1
  \end{cases}
\]
Rezultă că $f_n \xrar{PC} f$, unde funcția $f$ este definită de limita de mai sus.
Dar:
\begin{align*}
  ||f_n - f|| &= \sup_{x \in [0, 1)} |f_n(x) - f(x)| \\
              &= \max (\sup_{x \in [0, 1)} (|f_n(x) - f(x)|, |f_n(1) - f(1)|)) \\
              &= \max (\sup_{x \in [0, 1)} (x^n, 0)) \\
              &= 1,
\end{align*}
de unde obținem $\ds\lim_{n \to \infty} ||f_n - f|| = 1 \neq 0$. Rezultă că
$(f_n)$ este punctual convergent, nu uniform convergent pe $[0, 1)$.

\vspace{1cm}

Pentru caracterizarea convergenței, avem la dispoziție mai multe rezultate,
printre care și \emph{criteriul fundamental, al lui Cauchy}. Vom enunța, însă,
rezultatele care vor fi utile în special în rezolvarea exercițiilor:

\begin{theorem}[Integrare termen cu termen]\label{thm:uc-integrala}
  Fie $(f_n)$ un șir uniform convergent de funcții continue, $f_n : [a, b] \to \RR$.

  Atunci limita $f = \ds\lim_{n \to \infty} f_n$ este o funcție continuă pe
  $[a, b]$ și avem:
  \[
    \lim_{n \to \infty} \int_a^b f_n(x) dx = \int_a^b f(x) dx.
  \]

  Cu alte cuvinte, în cazul convergenței uniforme, limita comută cu integrala.
\end{theorem}

Pentru derivabilitate, avem:

\begin{theorem}[Derivare termen cu termen]\label{thm:uc-derivata}
  Fie $(f_n)$ un șir de funcții din $C^1([a, b])$ și $f, g$ funcții mărginite,
  cu $f, g : [a, b] \to \RR$. Dacă $f_n \xrar{PC} f$ și $f'_n \xrar{UC} g$ pe
  $[a, b]$, atunci $f$ este derivabilă pe $[a, b]$ și $f' = g$.
\end{theorem}

Pentru cazul șirurilor monotone de funcții continue, rezultatul următor arată
o legătură simplă între convergența punctuală și cea uniformă:

\begin{theorem}[U. Dini]\label{thm:dini}
  Fie $(f_n)$ un șir monoton de funcții continue, $f_n : [a, b] \to \RR$, astfel
  încît $f_n \xrar{PC} f$. Atunci $f_n \xrar{UC} f$.
\end{theorem}

Următorul rezultat va fi de folos în special în capitolul privitor la serii
de puteri:
\begin{theorem}[Stone---Weierstrass]\label{thm:label}
  Pentru orice funcție continuă $f : [a, b] \to \RR$, există un șir $(f_n)$
  \emph{de polinoame}, cu $f_n : [a, b] \to \RR$, astfel încît $f_n \xrar{UC} f$.
\end{theorem}

%%%%%%%%%%%%%%%%%%%%%%%%%%%%%%%%%%%%%%%%%%%%%%%%%%%%%%%%%%%%%%%%%%%%%%

\section{Serii de funcții}
\label{sec:serii-f}

Trecem acum la studiul seriilor de funcții, făcînd legături similare cu trecerea
de la șiruri de numere la serii de numere.

Începem cu o noțiune fundamentală:

\begin{definition}\label{def:pc-uc-serie-f}
  Mulțimea valorilor lui $x$ pentru care seria de funcții $\ds\sum_{n \geq 1} f_n(x)$
  este convergentă se numește \emph{mulțimea de convergență} a seriei, iar funcția
  $f : [a, b] \to \RR$, cu $f(x) = \ds\lim_{n \to \infty} S_n(x)$, unde $S_n(x)$
  este șirul sumelor parțiale pentru seria de funcții, se numește \emph{suma seriei}.

  
  Seria $\sum f_n$ se numește \emph{simplu (punctual) convergentă} către funcția $f$
  dacă șirul sumelor parțiale $(S_n(x))_n$ este punctual convergent către $f$.

  Similar, seria se numește \emph{uniform convergentă} către funcția $f$ dacă șirul
  sumelor parțiale este uniform convergent către $f$.

  Seria se numește \emph{absolut convergentă} dacă seria $\sum |f_n(x)|$ este simplu
  convergentă.
\end{definition}

Regăsim acum, atît noțiunile privitoare la șirurile de funcții, cît și criteriile
de convergență pentru serii de numere.

\begin{theorem}\label{thm:serie-f-integrala}
  Fie $\sum f_n$ o serie uniform convergentă de funcții continue, cu $f_n : [a, b] \to \RR$
  și fie $s$ suma acestei serii. Atunci $s$ este o funcție continuă pe $[a, b]$ și:
  \[
    \int_a^b s(x) dx = \sum_{n \geq 1} \int_a^b f_n(x) dx.
  \]
\end{theorem}

Avem rezultatul corespunzător și pentru derivate:

\begin{theorem}\label{thm:serie-f-derivata}
  Fie $\sum f_n$ o serie punctual convergentă de funcții din $C^1([a, b])$, cu suma $s$
  pe $[a, b]$ și astfel încît seria derivatelor $\sum f'_n$ să fie uniform convergentă.
  Atunci funcția $s$ este derivabilă pe $[a, b]$ și în plus:
  \[
    s'(x) = \sum_{n \geq 1} f'_n(x).
  \]
\end{theorem}

\begin{theorem}[Weierstrass]\label{thm:weierstrass-serii-f}
  Fie $\sum f_n$ o serie de funcții, cu $f_n : [a, b] \to \RR$ și fie $\sum a_n$ o serie
  convergentă de numere reale pozitive.

  Dacă $|f_n(x)| \leq a_n$, pentru orice $x \in [a, b]$ și $n \geq N$, cu $N$ fixat,
  atunci seria de funcții $\sum f_n$ este uniform convergentă pe $[a, b]$.
\end{theorem}

\begin{theorem}[Criteriul lui Abel]\label{thm:abel-serii-f}
  Dacă seria de funcții $\sum f_n$ se poate scrie $\sum \alpha_n v_n$, astfel încît
  seria de funcții $\sum v_n$ să fie uniform convergentă, iar $(\alpha_n)$ să fie un
  șir monoton de funcții egal mărginite (i.e.\ mărginite de aceeași constantă),
  atunci seria inițială este uniform convergentă.
\end{theorem}

O alternativă:

\begin{theorem}[Criteriul lui Dirichlet]\label{thm:dirichlet-serii-f}
  Dacă seria de funcții $\sum f_n$ se poate scrie sub forma $\sum \alpha_n v_n$,
  astfel încît șirul sumelor parțiale al seriei $\sum v_n$ să fie un șir de funcții
  egal mărginite, iar șirul $(\alpha_n)$ să fie un șir monoton de funcții, ce converge
  uniform către 0, atunci ea este uniform convergentă.
\end{theorem}

%%%%%%%%%%%%%%%%%%%%%%%%%%%%%%%%%%%%%%%%%%%%%%%%%%%%%%%%%%%%%%%%%%%%%%

\section{Formula lui Taylor}
\label{sec:taylor}

Putem asocia oricărei funcții cu anumite proprietăți \qq{bune} un polinom care o
aproximează. Este vorba despre \emph{polinomul Taylor}, definit astfel.

Fie $I \seq \RR$ un interval deschis și fie $f : I \to \RR$ o funcție de clasă
$C^m(I)$. Pentru orice $a \in I$, definim \emph{polinomul Taylor} de gradul $n \leq m$
asociat funcției $f$ în punctul $a$:
\[
  T_{n, f, a} = \sum_{k = 0}^n \frac{f^{(k)}(a)}{k!} (x - a)^k.
\]

Restul (eroarea aproximării) se definește prin:
\[
  R_{n, f, a}(x) = f(x) - T_{n, f, a} (x).
\]

Următorul rezultat ne arată că polinomul de mai sus poate fi transformat într-o formulă
mai exactă:

\begin{theorem}[Formula lui Taylor cu resturi Lagrange]\label{thm:taylor-lagrange}
  Fie $f : I \to \RR$ o funcție de clasă $C^{n+1}(I)$ și $a \in I$. Atunci, pentru
  orice $x \in I$, există $\xi \in (a, x)$ (sau $(x, a)$, după caz), astfel încît:
  \[
    f(x) = T_{n, f, a}(x) + \frac{(x - a)^{n+1}}{(n + 1)!} f^{n + 1}(\xi).
  \]
\end{theorem}

Privitor la restul rezultat din această formulă, avem următoarele:
\begin{itemize}
\item \emph{forma Peano}: $\exists \omega : I \to \RR$, cu
  $\ds\lim_{x \to a} \omega(x) = \omega(a) = 0$:
  \[
    R_{n, f, a}(x) = \frac{(x - a)^n}{n!} \omega(x).
  \]
\item \emph{forma integrală}:
  \[
    R_{n, f, a}(x) = \frac{1}{n!} \int_a^x f^{(n + 1)}(t) (x - t)^n dt.
  \]
\item $\ds\lim_{x \to a} \frac{R_{n, f, a}(x)}{(x - a)^n} = 0$.
\end{itemize}

%%%%%%%%%%%%%%%%%%%%%%%%%%%%%%%%%%%%%%%%%%%%%%%%%%%%%%%%%%%%%%%%%%%%%%

\section{Exerciții}
\label{sec:exercitii}

1. Să se studieze convergența simplă și uniformă a șirurilor de funcții:
\begin{enumerate}[(a)]
\item $f_n : [-1, 1] \to \RR, f_n(x) = \frac{x}{1 + nx^2}$;
\item $f_n : [0, 1] \to \RR, f_n(x) = x^n(1 - x^n)$;
\item $f_n : (-1, 1) \to \RR, f_n(x) = \frac{1 - x^n}{1 - x}$;
\item $f_n : [0, 1] \to \RR, f_n(x) = \frac{2nx}{1 + n^2 x^2}$;
\item $f_n : [0, \infty) \to \RR, f_n(x) = \frac{x + n}{x + n + 1}$;
\item $f_n : [0, \infty) \to \RR, f_n(x) = \frac{x}{1 + nx^2}$;
\item $f_n : \RR \to \RR, f_n(x) = \arctan (nx)$;
\item $f_n : \big[ 0, \frac{\pi}{2} \big] \to \RR, f_n(x) = n \sin^n x \cos x$.
\end{enumerate}

\vspace{1cm}

2. Să se arate că șirul $f_n : [0, \infty) \to \RR, f_n(x) = \frac{x}{x + n}$ nu converge
uniform pe $[0, \infty)$, dar converge uniform pe orice interval $[a, b]$, cu $0 < a < b$.

\vspace{1cm}

3. Să se arate că șirul de funcții $(f_n)$, unde:
\[
  f_n : (1, \infty) \to \RR, \quad f_n(x) = \sqrt{(n^2 + 1) \sin^2 \frac{\pi}{n} + nx} - \sqrt{nx}
\]
este uniform convergent (la 0).

\vspace{1cm}
4. Să se arate că șirul de funcții $(f_n)$, cu:
\[
  f_n : \RR \to \RR, \quad f_n(x) = \frac{1}{n} \arctan x^n
\]
converge uniform pe $\RR$, dar:
\[
  \Big( \lim_{n \to \infty} f_n(x) \Big)'_{x = 1} \neq \lim_{n \to \infty} f'_n(1).
\]

Rezultatele diferă deoarece șirul derivatelor nu converge uniform pe $\RR$.

\vspace{1cm}

5. Să se arate că șirul de funcții $(f_n)$, cu:
\[
  f_n : [0, 1] \to \RR, \quad f_n(x) = nxe^{-nx^2}
\]
converge, dar:
\[
  \lim_{n \to \infty} \int_0^1 f_n(x) dx \neq \int_0^1 \lim_{n \to \infty} f_n(x) dx.
\]

Rezultatul se explică prin faptul că șirul nu este uniform convergent. De exemplu,
pentru $x_n = \frac{1}{n} \in [0, 1]$, avem $f_n(x_n) \to 1$, dar în general,
$f_n(x) \to 0$.

\vspace{1cm}

6. Studiați convergența simplă și uniformă a șirului de funcții:
\[
  f_n : \Big[ 0, \frac{\pi}{2} \Big] \to \RR, \quad %
  f_n(x) = n \sin^n x \cos x.
\]



\end{document}